\usepackage[margin=3cm,bottom=2.5cm]{geometry}
\usepackage{hyperref}
\usepackage[usenames,dvipsnames]{color}
\usepackage{amssymb,amsmath,graphics,color,enumerate}
\usepackage{tikz}
\usepackage{xspace}
\usetikzlibrary{arrows,shapes,decorations,automata,backgrounds,petri,patterns} 
\usepackage{todonotes}
\usepackage[utf8x,utf8]{inputenc} 
\usepackage[T1]{fontenc}
\usepackage{cleveref}
%\usepackage[vlined,linesnumbered,boxruled]{algorithm2e}
%\usepackage[ruled]{algorithm2e}
%\renewcommand{\algorithmcfname}{Pseudocode}
\usepackage{booktabs}
\usepackage{array}
\usepackage{mathtools}
\usepackage{algorithm}
\usepackage{caption}
%\usepackage{dirtytalk}
\usepackage[noend]{algpseudocode}
\usepackage[numbers,sort]{natbib}
\captionsetup[algorithm]{
  labelfont = bf,
  labelsep = period
}

\usepackage{standalone}
\usepackage{tikz}
\usetikzlibrary{arrows}
\usetikzlibrary{patterns}

\usepackage{epsfig}
\usepackage{subfig}
%\usepackage{caption}
%\usepackage[T1]{fontenc}
\usepackage{amsmath,amsfonts,amssymb,epsfig,color,amsthm}
\usepackage{comment}
\usepackage{thmtools}
\usepackage{thm-restate}
\usepackage{mdframed}
\usepackage[absolute]{textpos}
\usepackage{enumitem}
\usepackage{soul}

\newtheorem{theorem}{Theorem}[section]
\newtheorem{definition}[theorem]{Definition}
\newtheorem{lemma}[theorem]{Lemma}
\newtheorem{corollary}[theorem]{Corollary}
\newtheorem{remark}[theorem]{Remark}
\newtheorem{proposition}[theorem]{Proposition}
\newtheorem{observation}[theorem]{Observation}
\newtheorem{question}[theorem]{Question}
\newtheorem{claim}{Claim}
\newtheorem{fact}{Fact}
\newtheorem{conjecture}{Conjecture}
\newtheorem{rtheorem}{Restated Theorem}
\newcommand{\retheorem}{theorem}

\newcommand{\Aa}{{\ensuremath{\mathcal{A}}}}
\newcommand{\Bb}{{\ensuremath{\mathcal{B}}}}
\newcommand{\Cc}{{\ensuremath{\mathcal{C}}}}
\newcommand{\Dd}{{\ensuremath{\mathcal{D}}}}
\newcommand{\Ee}{{\ensuremath{\mathcal{E}}}}
\newcommand{\Ff}{{\ensuremath{\mathcal{F}}}}
\newcommand{\Gg}{{\ensuremath{\mathcal{G}}}}
\newcommand{\Hh}{{\ensuremath{\mathcal{H}}}}
\newcommand{\Ii}{\mathcal{I}}
\newcommand{\Mm}{{\ensuremath{\mathcal{M}}}}
\newcommand{\Ss}{{\ensuremath{\mathcal{S}}}}
\newcommand{\Tt}{{\ensuremath{\mathcal{T}}}}
\newcommand{\Pp}{{\ensuremath{\mathcal{P}}}}
\newcommand{\Uu}{{\ensuremath{\mathcal{U}}}}
\newcommand{\Ww}{{\ensuremath{\mathcal{W}}}}
\newcommand{\nab}{\mathop{\triangledown}}
\newcommand{\cc}{\mathop{\rm{cc}}}
\newcommand{\poly}{\mathop{\rm{poly}}}
\newcommand{\Oh}{{\ensuremath{\mathcal{O}}}}
\newcommand{\Ohstar}{{\ensuremath{\mathcal{O}^*}}}

\newcommand{\PM}{{\ensuremath{\mathcal{PM}}}}

\newcommand{\Z}{\mathbb{Z}}
\newcommand{\F}{\mathbb{F}}
\newcommand{\Zq}{\mathbb{Z}_{\ge 0}}


\newcommand{\anonymyze}[1]{}%

\newcommand{\td}{\mathsf{td}}

\newcommand{\floor}[1]{\ensuremath{\left\lfloor{#1}\right\rfloor}}%
\newcommand{\ceil}[1]{\ensuremath{\left\lceil{#1}\right\rceil}}%
\newcommand{\ignore}[1]{}%
\newcommand{\field}[1]{\textup{GF}(#1)}

\newcommand{\ProblemFormat}[1]{\textsc{#1}}
\newcommand{\ProblemName}[1]{\ProblemFormat{#1}\xspace}

\newcommand{\MinHFunc}{\mathsf{MinHLT}}
\newcommand{\MaxHFunc}{\mathsf{MaxHLT}}
\newcommand{\MinHProb}{\ProblemName{MinHLT}}
\newcommand{\MaxHProb}{\ProblemName{MaxHLT}}
\newcommand{\DualMinHProb}{\ProblemName{DualMinHLT}}
\newcommand{\DualMaxHProb}{\ProblemName{DualMaxHLT}}
\newcommand\MinLLTProb{\ProblemName{MinLLT}}
\newcommand\MaxLLTProb{\ProblemName{MaxLLT}}


\newcommand{\Out}{\mathrm{Out}}
\newcommand{\In}{\mathrm{In}}
\newcommand{\First}{\mathrm{First}}
\newcommand{\Parent}{\mathrm{Parent}}
\newcommand{\Last}{\mathrm{Last}}
\newcommand{\Mult}{\mathrm{Mult}}
\newcommand{\Compr}{\mathrm{Compr}}

\DeclareMathOperator{\polylog}{polylog}


%%%%%% begin fancy problem definitions
\usepackage{mathrsfs}  
\usepackage{boxedminipage}
\usepackage{framed}
\usepackage{xifthen}
\usepackage{tabularx}
\usepackage{tikz} 
\usetikzlibrary{calc}

\newlength{\RoundedBoxWidth}
\newsavebox{\GrayRoundedBox}
\newenvironment{GrayBox}[1]%
   {\setlength{\RoundedBoxWidth}{.93\textwidth}
    \def\boxheading{#1}
    \begin{lrbox}{\GrayRoundedBox}
       \begin{minipage}{\RoundedBoxWidth}}%
   {   \end{minipage}
    \end{lrbox}
    \begin{center}
    \begin{tikzpicture}%
       \node(Text)[draw=black!20,fill=white,rounded corners,%
             inner sep=2ex,text width=\RoundedBoxWidth]%
             {\usebox{\GrayRoundedBox}};
        \coordinate(x) at (current bounding box.north west);
        \node [draw=white,rectangle,inner sep=3pt,anchor=north west,fill=white] 
        at ($(x)+(6pt,.75em)$) {\boxheading};
    \end{tikzpicture}
    \end{center}}     

\newenvironment{defproblemx}[2][]{\noindent\ignorespaces%
                                \FrameSep=6pt%
                                \parindent=0pt%
                \ifthenelse{\isempty{#1}}{%
                  \begin{GrayBox}{\textsc{#2}}%                
                }{%
                  \begin{GrayBox}{\textsc{#2} parameterized by~{#1}}%  
                }
                \newcommand\Prob{Problem:}%
                \newcommand\Input{Input:}%                        
                \begin{tabular*}{\textwidth}{@{\hspace{.1em}} >{\itshape} p{1.8cm} p{0.8\textwidth} @{}}%        
            }{
                \end{tabular*}%
                \end{GrayBox}%
                \ignorespacesafterend
            }  

\newcommand{\fancyparproblemdef}[4]{% FJR Version
  \begin{defproblemx}[#3]{#1}
    Input:  & #2 \\
    Question: & #4
  \end{defproblemx}
}%

\newcommand{\fancyproblemdef}[3]{% FJR Version
  \begin{defproblemx}{#1}
    Input:  & #2 \\
    Question: & #3
  \end{defproblemx}
}%

%%%%%%% END FANCY PROBLEM DEF


\usepackage{xspace}
\newcommand\mso[1]{\ensuremath{\mathsf{#1}}}
\newcommand\MSOone{\ensuremath{\mathsf{MSO}_1}\xspace}
\newcommand\MSOtwo{\ensuremath{\mathsf{MSO}_2}\xspace}

\newcommand\bags{\beta}
\newcommand\sig{\mathbf{s}}

\newcommand\tw{\mathsf{tw}}
\newcommand\const{\mathrm{DFS}}


\newcommand\rootmap{\mathfrak{r}}
\newcommand\roots{R}
\newcommand\heightmap{\mathfrak{h}}
\newcommand\aux{A\xspace}

\hyphenation{Tree-depth}

\newcommand\subparagraphstar[1]{\subparagraph*{#1}}
\newcommand\paragraphstar[1]{\paragraph*{#1}}
